\BeleskeID{03-s015}
% element: 03-s015:e01
% element: 03-s015:e02
% element: 03-s015:e03
% element: 03-s015:e04
% element: 03-s015:d01

U prikazanom uslovljenom MUX-u CS dozvoljava rad dekoderskog dela. Ako je \(Z_0\) izlaz neuslovljenog multipleksera, izlaz ove realizacije je \(Z=CS\,Z_0\). Kada je CS neaktivan, izlaz je logička nula, a ne visoka impedansa. Zato se izlazi takvih komponenti ne spajaju direktno samo na osnovu toga što je jedan CS aktivan.

Ako su \(g=2S_3+S_2\) i \(k=2S_1+S_0\), izlaz stabla 16/1 je \(Z=I_{4g+k}\). Na primer, adresa 1101 bira grupu 3 i lokalni ulaz 1, odnosno globalni \(I_{13}\). Svi početni multiplekseri istovremeno biraju po jedan podatak, ali završni propušta samo podatak iz adresirane grupe.

Lokalni nazivi I3--I0 ponavljaju se u svakom bloku, ali predstavljaju različite globalne podatke. Izlaz svakog početnog bloka dolazi na poseban ulaz završnog bloka; njihove izlazne žice se ne objedinjuju kratkim spojem. Ulaz CS prikazan je na osnovnoj komponenti levo; pri povezivanju stabla sve korišćene komponente moraju biti omogućene.
